% =============================================================================
% C_reproducibility.tex -- Appendix C: Reproducing numerical figures and builds
% Do NOT issue a second \appendix here -- main.tex already declares it once.
% =============================================================================

\section{Reproducing the numerical figures}
\label{app:reproducibility}

All plots in Secs.~\ref{sec:results} and~\ref{sec:discussion} are produced by
deterministic Python scripts in the repository directory \texttt{code/}, named
\texttt{calc-01.py} through \texttt{calc-10.py}.
Each script seeds its random number generators so that regenerated PDFs match
the versions checked into \texttt{manuscript/figures/} when the same defaults
are used.
The catalogue in \texttt{code/README.md} maps each script to its figure file and
manuscript section.

\subsection{Software requirements}

From the repository root, install Python~3.10 or newer and the pinned
dependencies:
\begin{verbatim}
python -m pip install -r code/requirements.txt
\end{verbatim}
A TeX distribution with \texttt{pdflatex} and \texttt{bibtex} is required only
if the full manuscript PDF is to be rebuilt (not for figure generation alone).

\subsection{Regenerating every figure}

The driver \texttt{scripts/build\_figures.py} executes every
\texttt{code/calc-*.py} in sorted order, writes \texttt{code/figures/figure\_N.pdf}
(and matching \texttt{.png} previews), and copies each PDF into
\texttt{manuscript/figures/} for \verb|\includegraphics|.
Equivalent entry points are \texttt{make figures} on POSIX systems and
\begin{verbatim}
python scripts/build_figures.py
\end{verbatim}
from the repository root on any platform.
To rebuild a single figure, pass the script stem, for example
\begin{verbatim}
python scripts/build_figures.py calc-03
\end{verbatim}

\subsection{Running one script interactively}

Each \texttt{calc-*.py} accepts a common command-line interface.
From the repository root, a typical invocation is
\begin{verbatim}
python code/calc-01.py --out code/figures/figure_1 --format pdf,png
\end{verbatim}
Optional flags include \texttt{--show} for an interactive Matplotlib window and
\texttt{--format} to restrict output types.
Script-specific options (for example \texttt{--gravity none|external|self} on
\texttt{calc-09.py}) are documented in the docstring at the top of each file.

\subsection{Physical Review figure geometry (\texttt{--prd-width})}

Before \verb|\includegraphics| scaling in the two-column \emph{Physical Review D}
layout, the Matplotlib canvas should match the intended float width.
Every figure script therefore accepts
\begin{verbatim}
--prd-width column|full
\end{verbatim}
where \texttt{column} sets the approximate single-column width in inches and
\texttt{full} sets the approximate two-column (\texttt{figure*}) width.
The manuscript build pipeline passes widths consistent with the float type of
each figure; when regenerating locally, choose \texttt{column} for
\texttt{figure} environments and \texttt{full} for \texttt{figure*} environments
so that font sizes and aspect ratios remain comparable to the published PDF.

\subsection{Building the manuscript PDF}

After figures are present under \texttt{manuscript/figures/}, compile from the
repository root with any of
\begin{verbatim}
make
pwsh scripts/build.ps1
bash scripts/build.sh
\end{verbatim}
The output is written to \texttt{build/main.pdf}.
Use \texttt{make figures} (or \texttt{python scripts/build\_figures.py}) before
the LaTeX step whenever a \texttt{calc-*.py} file has changed.
For text-only edits, \texttt{scripts/build.ps1 -SkipFigures} or
\texttt{bash scripts/build.sh --skip-figures} skips the numerics pass.
